\subsection{The actual three-torsion grid and its completed normalizer}

\begin{result}{E.184}{Proved}
{Positive fusion and full normalizer recognition for the labelled grid}
Let \(T=\mathbb R^2/\mathbb Z^2\) be an oriented torus, mark
\(T[3]\), label its origin by zero, and put
\[
 \mathcal M=\Gamma_{1,[9]},\qquad
 H=T[3]\cong(\mathbb Z/3)^2\leq\mathcal M,
\tag{E.184.1}
\]
where \(H\) acts by translations, with its specified regular
label permutation.  The normalized full-GT action of E.182
can be conjugated by one element of the completed pure
positive mapping class group so as to fix \(H\) pointwise.

The actual ordinary normalizer and its zero-label stabilizer are
\[
 \begin{gathered}
 \mathcal N=N_{\mathcal M}(H)
       =H\rtimes\operatorname{SL}_2(\mathbb Z),\qquad
 \mathcal N_0\cong\operatorname{SL}_2(\mathbb Z),\\
 \mathcal N_0\xrightarrow[\text{forget the other eight marks}]{\sim}
                         \Gamma_{1,1}.
 \end{gathered}
\tag{E.184.2}
\]
They have their full induced profinite topologies, and
\[
                 N_{\widehat{\mathcal M}}(H)
                         =\widehat{\mathcal N}.
\tag{E.184.3}
\]
Thus the corrected action restricts to this actual completed
normalizer, and its action on \(\widehat{\mathcal N}_0\),
under E.184.2, has the original one-point outer action.
These assertions establish the qualitative translation
normalization in the historical CM comparison.  They do
not identify a higher-point elliptic-involution quotient
with a genus-zero action, or certify the old finite bounds.
\end{result}

\begin{proof}
We retain the surface representation supplied by E.182.
Write \(\Pi=\pi_1(T\setminus T[3])\), a free group of rank ten.
The completed mixed mapping class group acts on
\(\widehat\Pi\) by outer automorphisms.  The automorphism
\(\alpha_f\) of \(\widehat{\mathcal M}\) is compatible with
an automorphism \(\eta_f\) of \(\widehat\Pi\), which preserves
each positive peripheral conjugacy class.  Killing eight
of the nine peripheral classes gives the once-punctured
torus free group, with its last peripheral class a
commutator.  In its class-two quotient the commutator
exponent is the determinant on compact homology.
It follows that the compact matrix of \(\eta_f\) has
determinant one.  In particular \(\alpha_f\) preserves
the kernel of the ordinary compact-period map modulo six:
\[
 \operatorname{per}_6:\mathcal M
              \longrightarrow\operatorname{SL}_2(\mathbb Z/6).
\tag{E.184.4}
\]
Indeed its effect on that matrix representation is
conjugation by the compact matrix of \(\eta_f\).
Its effect on the separate label quotient is the identity.

Let \(P\) be the subgroup that fixes every label and has
compact period in \(\Gamma(6)\).  Let
\(E=P H=P\rtimes H\), an actual finite-index subgroup of
\(\mathcal M\).  Both \(\widehat P\) and \(\widehat E\)
are preserved by \(\alpha_f\), by E.184.4 and the label
condition; the quotient action on \(E/P=H\) is the identity.

Here is the geometric space and finite-cover argument needed
to turn that last statement into fusion.  Over the modular
curve of full level six, let \(X\) be the space of eight
ordered, distinct, nonzero points on the universal torus;
the zero point is the ninth mark.  Equivalently use the
upper half-plane with these configurations and quotient by
\(\Gamma(6)\).  The base is a smooth aspherical noncompact
curve, since \(\Gamma(6)\) is torsion free and free of
finite rank.  Successively forgetting the eight points
has punctured-torus fibres, with free ranks
\(2,3,\ldots,9\).  These are locally trivial configuration
fibrations.  Thus \(X\) is a smooth aspherical manifold
of real dimension eighteen and \(\pi_1(X)=P\) is an
iterated extension of finitely generated free groups of
rank at least two.

Index the labels by \(h\in H\), and write a normalized
configuration as \((p_h)\), with \(p_0=0\).
The actual finite action of \(k\in H\) on \(X\) is
\[
                   (p_h)_h\longmapsto(p_{h+k}-p_k)_h.
\tag{E.184.5}
\]
The base torsion grid is fixed, and its orbifold path
extension is precisely \(E=P\rtimes H\).

The proof of E.176 applies to this \(C_3^2\)-action after
the following explicit changes.  Use \(\mathbb F_3\)
throughout.  For a free group \(F\) of rank \(r\geq2\),
put \(N=F'F^3\) and \(W=F/(N'N^3)\).
In the regular \(C_3^r\)-cover of the rose, the cube of
the \(j\)-th generator is the nonzero cycle on the three
\(j\)-edges supported on \(\langle e_j\rangle\).
For any nonzero translation choose \(j\) whose line does
not contain it.  This moves that cycle.  Hence the deck
action on \(H_1(N,\mathbb F_3)\) is faithful, and the
center of the finite group \(W\) projects trivially to
\(F/N\).  The group has order
\(3^{r+1+3^r(r-1)}\).
The same two-row Lyndon--Hochschild--Serre induction,
with its required two iterations, therefore kills every
positive mod-three cohomology restriction after a finite
characteristic refinement.

For the cyclic fixed-point pass replace the two maps in
the complete resolution by
\[
                    1-a,\qquad 1+a+a^2.
\tag{E.184.6}
\]
On a fixed cochain both maps are zero over \(\mathbb F_3\).
On free-orbit product cochains the Tate complex is exact,
by the same coinduction argument as in E.176.
The vertical cochain range remains bounded; the reduced
cone, including its degree \(-1\) when a fixed set is
empty, is unchanged.  Applying this pass first to one
\(C_3\), then to its quotient in \(H\), proves both
the exact finite centralizer projection and the simultaneous
finite fusion assertion of E.176.2--E.176.4 for this space.
The proof still uses actual finite characteristic covers
and fixed paths, not an assertion that a discrete fixed
locus is already a profinite centralizer.

The fixed components can be determined completely.
Equation E.184.5 is fixed for every \(k\) exactly when
\[
                         p_{h+k}=p_h+p_k.
\tag{E.184.7}
\]
The distinctness condition says that \(h\mapsto p_h\)
is an isomorphism \(H\to T[3]\).
The level-six marking trivializes \(T[3]\), so the
fixed locus has exactly forty-eight components, each
a section over the same connected modular curve:
\[
                          X^H=\coprod_{D\in\mathrm{GL}_2(\mathbb F_3)}X_D.
\tag{E.184.8}
\]
Use one reference fibre and paths to those forty-eight
points to obtain the ordinary complements \(H_D\leq E\).
For each finite quotient, finite fusion expresses the
image of \(\alpha_f(H)\) as a \(P\)-conjugate of some
\(H_D\), with the parameterization by the named label
group \(H\) retained pointwise.
The finite set of components and the compact group
\(\widehat P\) give a finite-intersection argument:
the closed solution sets in
\(\widehat P\times\mathrm{GL}_2(\mathbb F_3)\)
for all finite quotients are nonempty and have the
finite-intersection property.  Therefore
\[
                    \alpha_f(H)=uH_Du^{-1}
                  \quad\text{pointwise},\qquad u\in\widehat P,
\tag{E.184.9}
\]
for one component \(D\).

We now distinguish the two orientations, rather than
silently discard twenty-four fixed components.
Choose \(A_D\in\mathrm{GL}_2(\mathbb Z)\) reducing to
\(D\), with determinant \(\varepsilon\in\{1,-1\}\).
Such a lift exists by elementary row operations modulo
three.  The corresponding linear homeomorphism sends
the standard grid \(p_h\) to \(p_{Dh}\).
Choose an identity-isotopic orientation-preserving
homeomorphism \(k_D\) moving these labelled points to
the same grid, as supplied by the selected configuration
path.  Then
\[
                         g_D=A_D^{-1}k_D
\tag{E.184.10}
\]
fixes each original label and has orientation
\(\varepsilon\).  Transport of the deck translations
along that path gives, pointwise,
\(H_D=g_D^{-1}Hg_D\).
Thus E.184.9 supplies a pure, possibly orientation-reversing
conjugator \(h=ug_D^{-1}\).  The corrected surface
automorphism has peripheral exponent \(\varepsilon\)
and centralizes the outer action of each named element
of \(H\).

There is an exact integral test for its sign.  The quotient
of the punctured torus by the free translation action is
the punctured torus with one puncture.  Its actual group
extension is
\[
 1\longrightarrow\widehat\Pi
 \longrightarrow\widehat F(A,B)
 \longrightarrow H\longrightarrow1,
 \qquad H=F(A,B)^{\rm ab}/3,
\tag{E.184.11}
\]
and its peripheral commutator \(c\) belongs to \(\Pi\).
All nine punctures above it have ramification index one.
The conjugation action embeds the middle group into
\(\operatorname{Aut}(\widehat\Pi)\): its kernel intersects
the open center-free subgroup \(\widehat\Pi\) trivially,
so is finite, whereas a free profinite group has no
nontrivial finite subgroup.  Its image is exactly the
inverse image of the outer deck group \(H\), since it
contains \(\operatorname{Inn}(\widehat\Pi)\) and maps
onto \(H\).
Therefore any surface automorphism normalizing this
outer deck group extends, by conjugation, to an
automorphism of the actual group in E.184.11.

For the corrected automorphism this extension acts
identically on \(H\).  It takes \(c\) to a conjugate
of \(c^{\varepsilon}\), because it has that exponent
on every puncture above \(c\).
The determinant of its matrix on
\(\widehat F(A,B)^{\rm ab}=\widehat{\mathbb Z}^2\)
is consequently \(\varepsilon\), as seen in the
class-two commutator coordinate.  The same matrix is
the identity modulo three, by the action on \(H\).
Hence
\[
                            \varepsilon\equiv1\pmod3,
                        \qquad \varepsilon=1.
\tag{E.184.12}
\]
The correcting \(g_D\) and \(h\) are thus positive.
This proves the first assertion, including purity and
pointwise rather than unparameterized fusion.

We give the normalizer comparison as well, since its
completion is a separate issue.  In the ordinary category
the punctured Dehn--Nielsen--Baer theorem makes
\(\mathcal M\to\operatorname{Out}(\Pi)\) injective:
the surface has genus one and nine punctures, hence
negative Euler characteristic.  The exact punctured
statement and proof are
\citet[version 5.0, \S8.2.7, Theorem 8.8,
pp.245--246]{LitFarbMargalit2011Version5}.
We use this ordinary faithfulness in the kernel calculation
below; no profinite faithfulness is asserted here.
In the ordinary category
the same deck extension embeds \(F(A,B)\) in
\(\operatorname{Aut}(\Pi)\).  A normalizing mapping class
extends to a peripheral-preserving outer automorphism of
\(F(A,B)\).  The rank-two Nielsen identification
\(\operatorname{Out}(F_2)=\mathrm{GL}_2(\mathbb Z)\),
with the peripheral sign given by determinant, identifies
the positive quotient with \(\mathrm{SL}_2(\mathbb Z)\);
see also the explicit exact sequence in
\citet[\S3, p.335]{LitBuxErshovRapinchuk2011}.
The kernel is exactly the deck group \(H\): an inner
automorphism of \(F(A,B)\) restricts, modulo inner
automorphisms of \(\Pi\), to its deck class; subtracting
that actual deck transformation leaves a mapping class
with trivial outer action on \(\Pi\), which is the
identity by the stated injectivity.
Conversely every integral positive linear map of the
torus normalizes \(T[3]\).  It lifts the same quotient
map and gives a splitting.  Exactly one deck translate
of a lift fixes zero.  This proves E.184.2.
The same compact lattice is carried through quotienting
by \(H\) and through forgetting the other points: the
quotient is the multiplication-by-three map, whose
compact homology matrix is \(3I\).  Thus no different
modular marking is used in the two maps.

Every finite quotient of \(\mathcal N_0\) extends to
the zero-label stabilizer in \(\mathcal M\), by the
ordinary forgetting map in E.184.2.
That stabilizer has index nine in \(\mathcal M\), so
its finite quotients are retained in finite quotients
of \(\mathcal M\), for example by the ordinary induced
coset representation.  This proves the full induced
topology of \(\mathcal N_0\).
Joining the label map retains the other eight cosets
in \(\mathcal N\), proving its full topology as well.

Finally take \(v\in N_{\widehat{\mathcal M}}(H)\).
Its permutation of the nine labels lies in the affine
normalizer of the regular translation group, so has
the form \((t,D)\) with
\(t\in H\) and \(D\in\mathrm{GL}_2(\mathbb F_3)\).
Apply the extension argument E.184.11 to the outer
surface action of \(v\).  Its peripheral exponent is
one.  Its matrix on the quotient-torus compact lattice
has determinant one and reduces to \(D\) modulo three.
It is also the compact matrix \(\operatorname{per}(v)\)
on the covering torus: both matrices commute with the
covering map \(3I\), and multiplication by three is
injective on \(\widehat{\mathbb Z}\).
In particular
\[
                   D=\operatorname{per}(v)\bmod3
                             \in\mathrm{SL}_2(\mathbb F_3).
\tag{E.184.13}
\]
Choose \(A\in\mathrm{SL}_2(\mathbb Z)\) with
\(A\equiv\operatorname{per}(v)\pmod6\), and put
\(n=t\,s_{\rm lin}(A)\in\mathcal N\).
The surjectivity of reduction modulo six follows by
elementary matrices and the Chinese remainder theorem.
The elements \(n,v\) have the same label permutation,
the same compact period modulo six, and the same
pointwise conjugation on \(H\).
Consequently \(n^{-1}v\in C_{\widehat P}(H)\).
The mod-three finite centralizer statement proved
above, followed by the compact inverse limit, gives
\[
                         C_{\widehat P}(H)
                                      =\overline{C_P(H)}.
\tag{E.184.14}
\]
Thus \(v\in\overline{\mathcal N}\), proving E.184.3;
the reverse inclusion is immediate.

An automorphism fixing \(H\) pointwise preserves this
normalizer.  The correcting conjugator is pure, so it
also preserves the zero-label condition.
Ordinary forgetting, already proved in E.182, identifies
the restriction on \(\widehat{\mathcal N}_0\) with the
original one-point action up to the image of that same
conjugator.  This is one completed comparison, not
separate choices after different coefficient reductions.
The elliptic-involution and subsequent branch maps still
require their own comparison.
\end{proof}
